\section{Related Work}
\label{sec:related}

\noindent\textbf{Device fingerprinting.} Devices are identifiable from wire-side behavior without
payload access: from remote clock skew~\cite{kohno2005}, from timing signatures that reveal device and
device type~\cite{gtid2015}, and, for cyber-physical systems, from cross-layer response timing,
physical-operation timing, and device physics~\cite{formby2016,gu2018}. This line defines the threat
we address rather than a defense. Our objective differs: we suppress the master-facing CLRT and
segment-shape features these techniques exploit, for one evaluated device and its transaction classes.

\noindent\textbf{Network-timing traffic analysis and its defenses.} Inter-event timing leaks content,
from keystroke intervals on SSH~\cite{song2001} to timing-only web-page recovery~\cite{feghhi2016} and
deep-learning website fingerprinting~\cite{sirinam2018}, though such accuracy is contested at realistic
open-world scale~\cite{panchenko2016}. Defenses shape the traffic: adaptive padding breaks the
inter-packet pattern~\cite{juarez2016}, comprehensive side-channel mitigation shapes a tenant's
traffic to a leak-free schedule~\cite{pacer2022}, and anonymity systems resist traffic analysis at the
network layer~\cite{hornet2015,taranet2018}. These target browser or cloud traffic and assume a
cooperating end host or hypervisor that can add cover traffic. Our defense normalizes the
acknowledgment-to-response timing of a protocol-bound DNP3 exchange in the network, with no endpoint
change and no added application bytes, and fixes one cross-layer interval rather than obscuring a
distribution.

\noindent\textbf{Packet-size padding, morphing, and constant-rate shaping.} Size defenses reshape a
flow's length signature: morphing transforms one class's distribution into another's~\cite{wright2009},
dependent link padding bounds low-latency leakage~\cite{wang2008dlp}, and constant-rate and
burst-molding schemes fix the rate at a cost that is acceptable for one flow but prohibitive across a
device fleet~\cite{tamaraw2014,walkietalkie2017}; later analysis shows per-packet padding and
fragmentation still leave exploitable features~\cite{dyer2012}. The same size-and-rate channel
identifies devices even under encryption, where shaping, made differentially private, is the
countermeasure~\cite{apthorpe2019,netshaper2024}. These change or pad bytes and target the end host.
Our mechanism neither pads nor rewrites DNP3 bytes; it re-segments the response only at existing
CRC-block boundaries, so the reassembled message and every per-block CRC are preserved while the
observed shape is fixed.

\noindent\textbf{Reconnaissance, moving-target defense, and DNP3 security.} A body of work disrupts
grid reconnaissance rather than detecting its use: DefRec virtualizes physical functions to present
deceptive content and connectivity~\cite{defrec2020}, RAINCOAT randomizes data acquisition and injects
decoys~\cite{raincoat2019}, and a companion testbed evaluates such approaches on grid cyber-physical
infrastructure~\cite{lintestbed2020}; the reconnaissance these blunt ultimately enables
false-data-injection and process-control attacks~\cite{liufdia2009,cardenas2011,sridhar2012}. A
parallel line detects rather than obfuscates, adapting specification-based and state-based intrusion
detection to DNP3 and analyzing its attack surface~\cite{linbro2013,fovino2010,east2009}, and in-network
obfuscation has denied reconnaissance by rewriting data-plane topology against link-flooding
attacks~\cite{nethide2018}. These operate at the content, connectivity, topology, or detection layer.
Our defense instead obfuscates the wire-visible timing and shape of the transport exchange itself,
below the semantics those defenses reshape, and complements rather than replaces them. Standards and
incident analyses frame the setting~\cite{dnp3std,ukraine2015,nist80082}.

\noindent\textbf{Timing and shaping on programmable switches.} Programmable data planes make in-network
shaping practical. The P4 model exposes a programmable parser and match-action
pipeline~\cite{p4ccr2014}; push-in-first-out queues define ideal programmable
scheduling~\cite{pifo2016}, and SP-PIFO approximates it with a few strict-priority queues on today's
switches~\cite{sppifo2020}, the primitive class our deadline scheduler uses. Recent systems bring
traffic-analysis defense fully into the data plane: Ditto pads and injects chaff to a fixed WAN
pattern~\cite{ditto2022}, Securitas obfuscates flows with flexible in-network
shaping~\cite{securitas2026}, Minos mounts a dynamic defense on programmable switches~\cite{minos2025},
and, closest to our setting, Hu and Lin enhance industrial-protocol security with programmable
switches~\cite{hulin2023}. Our defense differs from these in two ways: it is byte-preserving and carves
only at existing DNP3 boundaries, so the outstation's protocol remains valid, and it adds a
control-command treatment that hides operation timing from the master-facing observer without exposing
the hold.

\section{Conclusion}
\label{sec:concl}

An observer of unencrypted DNP3 traffic can fingerprint an outstation from the cross-layer response
time and the segment shape of its responses, and neither encryption, endpoint changes, nor generic
traffic-analysis defenses fit the constraints of a deployed DNP3 link. We placed the defense in the
data plane of one programmable switch, where it holds the acknowledgment and response and releases them
at fixed offsets from the request to normalize the response time, replicates and carves the eligible
response into a fixed $[28,21]$ shape at an existing CRC boundary, and holds a control command to pin
its master-visible timing to the request. On a physical SEL-751 relay the defended cross-layer response
time is a fixed 4.001~ms, every eligible response presents the $[28,21]$ shape with a CRC- and
checksum-valid reconstruction, the control-command gap stays at 4.00~ms across the evaluated internal
delays, and a read-versus-select timing classifier falls to chance, within a single-device,
master-facing evaluation boundary whose relay-facing verification we leave to future work.
